\section{Introduction}

Vision-language-action (VLA) policies make language-conditioned robot
manipulation increasingly general, but reliable control requires spatial
grounding in the robot's own embodiment. Recent robot foundation models show
that broad pretraining and large robot datasets can yield policies that transfer
across tasks and embodiments~\citep{brohan2022rt1,brohan2023rt2,team2024octo,kim2024openvla}.
For manipulation, however, recognizing the instructed object is not sufficient:
the policy must relate visual evidence to the robot's proprioceptive state and
to the geometry of the surrounding scene. Spatial robustness under changes in
visual perspective is therefore a core requirement for dependable VLA control.

A key unresolved failure mode is that a VLA policy can be task-grounded without
being self-grounded. In controlled camera-perturbation evaluations where the
underlying robot-scene state is unchanged, lateral camera shifts induce
co-directional shifts in predicted robot trajectories. This systematic behavior
suggests that current VLAs may exploit image-space shortcuts, such as
pixel-coordinate correlations, instead of anchoring actions to proprioceptive
state and robot-object geometry. A policy trained only to match demonstrations
under a fixed observation distribution can therefore appear competent while
remaining brittle when the same physical state is viewed from a different
perspective.

Existing grounding-oriented VLA methods show that perception, attention
allocation, and embodiment cues must be treated as first-class control
problems. Recent methods improve spatial awareness with visual trajectory
prompts, 3D representations, camera-space action formulations, view synthesis,
or explicit treatments of proprioception~\citep{zheng2025tracevla,qu2025spatialvla,zhang2025ocvla,gu2026vistabot,zhao2025statefree,lu2026gap}.
These approaches are valuable, but they do not directly provide a
self-supervised signal that tells the policy where its own acting body appears
in the image. For spatially reliable manipulation, the missing cue is not only
object grounding or view robustness, but a binding between proprioceptive state
and visual evidence of the robot itself.

\textsc{ProprioVLA} addresses this missing-self problem by learning
self-supervised proprioceptive grounding from robot interaction. The key
insight is that the robot reveals itself through visual motion synchronized
with changes in proprioceptive state. Rather than using proprioception as a
passive input token, \textsc{ProprioVLA} uses proprioceptive state to query
visual features and derives visual supervision from state-aligned residual
flow. This signal highlights residual motion that is
both locally salient in the image and aligned with proprioceptive change,
indicating where the robot's own visual evidence appears during interaction.

The method realizes this idea through confidence-weighted self-grounding
heatmaps and two complementary objectives. First, \textsc{ProprioVLA} mines
heatmaps from robot interaction by selecting residual motion that is salient in
the image and aligned with proprioceptive change, while downweighting uncertain
frames. These heatmaps supervise a visual self-grounding objective that aligns
state-conditioned attention with robot-related visual evidence. Second, a
proprioceptive reconstruction objective requires the attended visual
representation to recover the robot's own state, preventing self-grounding from
collapsing into a purely spatial attention cue. Together, these objectives
provide proprioceptive grounding without manual masks, external grounding
datasets, or supervision from stronger foundation models.

Experiments test whether proprioceptive self-grounding improves spatial
reliability beyond ordinary action supervision and state concatenation. We
evaluate \textsc{ProprioVLA} in real-robot and simulation environments,
including controlled camera-perturbation settings that directly probe view
robustness. Across these evaluations, \textsc{ProprioVLA} improves task success
and robustness over action-supervised VLA adaptation. Ablations further support
the role of both components: state-aligned heatmaps provide spatial supervision
for proprioceptive visual queries, while reconstruction anchors the attended
representation to the robot state. These results suggest that self-supervised
proprioceptive grounding is a practical route toward spatially reliable VLA
control.

Our contributions are:
\begin{itemize}
    \item We identify a systematic grounding failure in VLA policies: camera
    shifts can induce co-directional shifts in predicted robot trajectories,
    revealing reliance on image-space shortcuts rather than proprioceptive
    grounding.
    \item We introduce state-aligned residual flow as an implicit supervision
    signal for self-grounding, enabling confidence-weighted heatmaps to be mined
    from robot interaction without manual annotations, external grounding
    datasets, or stronger foundation-model supervision.
    \item We propose \textsc{ProprioVLA}, a self-supervised learning framework
    that couples state-conditioned visual attention with proprioceptive
    reconstruction to improve task success and robustness in real-robot and
    simulation environments.
\end{itemize}
